\documentclass[a4paper,12pt]{article}

% ------------------------------------------------
% Кодировка, язык, математика
% ------------------------------------------------
\usepackage[T2A]{fontenc}
\usepackage[utf8]{inputenc}
\usepackage[russian]{babel}
\usepackage{amsmath,amssymb,amsthm,mathtools}
\usepackage{icomma}

% ------------------------------------------------
% Шрифт и типографика
% ------------------------------------------------
\usepackage{helvet}
\renewcommand{\familydefault}{\sfdefault}
\usepackage{microtype}
\usepackage{setspace}
\onehalfspacing
\usepackage{parskip}

% ------------------------------------------------
% Страница
% ------------------------------------------------
\usepackage[
  a4paper,
  left=20mm,
  right=20mm,
  top=20mm,
  bottom=22mm
]{geometry}

% ------------------------------------------------
% Графика
% ------------------------------------------------
\usepackage{tikz}
\usetikzlibrary{calc}
\usetikzlibrary{arrows.meta,positioning,shapes.geometric}

\usepackage{tkz-euclide}

\usepackage{pgfplots}
\pgfplotsset{compat=1.18}

% ------------------------------------------------
% Таблицы и списки
% ------------------------------------------------
\usepackage{tabularx,booktabs,longtable}
\usepackage{enumitem}

% ------------------------------------------------
% Служебные пакеты
% ------------------------------------------------
\usepackage{xcolor}
\usepackage{hyperref}
\usepackage{fancyhdr}

% ------------------------------------------------
% Палитра
% ------------------------------------------------
\definecolor{accent}{HTML}{008080}
\definecolor{darkaccent}{HTML}{004D4D}
\definecolor{textgray}{HTML}{333333}
\definecolor{softgray}{HTML}{E9EEEE}
\definecolor{gridgray}{HTML}{D6DEDE}

\color{textgray}

% ------------------------------------------------
% Ссылки
% ------------------------------------------------
\hypersetup{
  colorlinks=true,
  linkcolor=darkaccent,
  urlcolor=darkaccent,
  pdfborder={0 0 0}
}

% ------------------------------------------------
% Колонтитулы
% ------------------------------------------------
\pagestyle{fancy}
\fancyhf{}
\fancyhead[L]{\small\color{darkaccent}Физика}
\fancyhead[R]{\small\color{textgray}График координаты \(x(t)\)}
\fancyfoot[C]{\small\thepage}
\renewcommand{\headrulewidth}{0.3pt}
\renewcommand{\footrulewidth}{0pt}

\fancypagestyle{plain}{
  \fancyhf{}
  \renewcommand{\headrulewidth}{0pt}
  \renewcommand{\footrulewidth}{0pt}
}

% ------------------------------------------------
% Заголовки
% ------------------------------------------------
\usepackage{titlesec}

\titleformat{\section}
  {\Large\bfseries\color{darkaccent}}
  {}
  {0pt}
  {}

\titleformat{\subsection}
  {\large\bfseries\color{accent}}
  {}
  {0pt}
  {}

\titlespacing*{\section}{0pt}{1.2em}{0.6em}
\titlespacing*{\subsection}{0pt}{1em}{0.4em}

% ------------------------------------------------
% Настройки списков
% ------------------------------------------------
\setlist[itemize]{
  leftmargin=1.6em,
  itemsep=0.25em,
  topsep=0.3em
}

\setlist[enumerate]{
  leftmargin=1.8em,
  itemsep=0.35em,
  topsep=0.4em
}

% ------------------------------------------------
% TikZ: стили блок-схем
% ------------------------------------------------
\tikzset{
  flowstart/.style={
    ellipse,
    draw=accent,
    line width=0.9pt,
    minimum width=4.6cm,
    minimum height=1cm,
    align=center,
    inner sep=6pt
  },
  flowaction/.style={
    rounded corners=4pt,
    rectangle,
    draw=accent,
    line width=0.9pt,
    text width=10.8cm,
    minimum height=1cm,
    align=center,
    inner sep=7pt
  },
  flowdecision/.style={
    diamond,
    aspect=2.7,
    draw=accent,
    line width=0.9pt,
    text width=5.2cm,
    align=center,
    inner sep=2pt
  },
  flowarrow/.style={
    -{Latex[length=3mm,width=2mm]},
    line width=0.9pt,
    draw=darkaccent
  }
}

% ------------------------------------------------
% Удобные команды
% ------------------------------------------------
\newcommand{\key}[1]{\textbf{\color{darkaccent}#1}}
\newcommand{\important}[1]{\textbf{\color{accent}#1}}

\begin{document}

% ================================================================
% ТИТУЛЬНЫЙ ЛИСТ
% ================================================================
\begin{titlepage}
\thispagestyle{plain}

\begin{center}

\vspace*{2.5cm}

{\Large\color{accent}\textbf{ФИЗИКА}}

\vspace{1.6cm}

{\fontsize{28}{34}\selectfont
\bfseries
Чек-лист}

\vspace{0.7cm}

{\LARGE
График зависимости координаты\\[0.2cm]
от времени \(x(t)\)}

\vspace{1cm}

{\large
Координата, направление движения\\
и скорость тела}

\vfill

{\large
Кирилл Зиновьев, 7 класс}

\vspace{0.5cm}

{\normalsize
\today}

\vspace{1.2cm}

{\small
Лёвин Артём Александрович\\
эксклюзивно для Кирилла Зиновьева}

\vspace{1.8cm}

\begin{tikzpicture}[remember picture,overlay]
  \draw[
    accent,
    line width=1.2pt
  ]
  ([xshift=25mm,yshift=17mm]current page.south west)
    .. controls
  ([xshift=65mm,yshift=27mm]current page.south west)
    and
  ([xshift=135mm,yshift=7mm]current page.south west)
    ..
  ([xshift=-25mm,yshift=20mm]current page.south east);
\end{tikzpicture}

\end{center}

\end{titlepage}

% ================================================================
% ОСНОВНОЙ ЧЕК-ЛИСТ
% ================================================================
\section{1. Что показывает график \(x(t)\)}

На графике зависимости координаты от времени:

\[
\boxed{\text{по горизонтали откладывается время }t}
\]

\[
\boxed{\text{по вертикали откладывается координата }x}
\]

Каждая точка графика отвечает на вопрос:

\begin{center}
\textit{«Где находилось тело в данный момент времени?»}
\end{center}

Например, если точка графика имеет координаты

\[
t=2~\text{с},
\qquad
x=3~\text{м},
\]

то через \(2\) секунды после начала наблюдения тело находилось в точке
с координатой \(3\) м.

\subsection{Ключевые точки графика}

При первом взгляде на график полезно сразу отметить:

\begin{itemize}
  \item \key{начальную точку};
  \item \key{точки пересечения графика с осью времени};
  \item \key{точки излома};
  \item конечную точку рассматриваемого участка.
\end{itemize}

Точка пересечения с горизонтальной осью особенно важна:

\[
x=0.
\]

Это означает, что в данный момент тело находится в \important{начале координат}.

\subsection{Пример}

\begin{center}
\begin{tikzpicture}
\begin{axis}[
  width=12.5cm,
  height=7cm,
  axis lines=middle,
  xmin=-0.3,
  xmax=8.5,
  ymin=-5,
  ymax=5,
  xlabel={$t,\ \text{с}$},
  ylabel={$x,\ \text{м}$},
  xtick={0,1,2,3,4,5,6,7,8},
  ytick={-4,-3,-2,-1,0,1,2,3,4},
  grid=both,
  major grid style={draw=gridgray},
  minor tick num=0,
  clip=false,
  unit vector ratio=1 1
]
\addplot[
  accent,
  very thick,
  mark=*,
  mark size=2pt
]
coordinates {
  (0,4)
  (3,-4)
  (8,2)
};

\node[above right] at (axis cs:0,4) {$A$};
\node[below] at (axis cs:3,-4) {$C$};
\node[above] at (axis cs:8,2) {$E$};

\end{axis}
\end{tikzpicture}
\end{center}

На участке \(AC\) время увеличивается, а координата уменьшается.

Следовательно, тело движется в сторону \important{уменьшения координаты \(x\)}.

В точке \(C\) линия меняет наклон. После неё координата начинает
увеличиваться. Значит, тело \important{изменило направление движения}.

% ================================================================
\section{2. Как определить направление движения}

Главный ориентир --- изменение координаты.

\begin{center}
\begin{tabularx}{0.92\linewidth}{@{}X X@{}}
\toprule
\textbf{Что происходит с \(x\)} &
\textbf{Что происходит с телом} \\
\midrule
\(x\) увеличивается &
тело движется вдоль положительного направления оси \(Ox\) \\
\(x\) уменьшается &
тело движется против положительного направления оси \(Ox\) \\
\(x\) постоянно &
тело покоится \\
\bottomrule
\end{tabularx}
\end{center}

\subsection{Очень важное различие}

Если график пересекает горизонтальную ось,

\[
x=0,
\]

это означает только то, что тело прошло через начало координат.

Из равенства

\[
x=0
\]

\important{не следует}, что

\[
v=0.
\]

Тело в этот момент может продолжать двигаться.

\subsection{Как распознать разворот}

Если до некоторого момента координата уменьшалась, а после него
увеличивается, тело в этот момент меняет направление движения.

Схематично:

\[
\text{координата уменьшается}
\longrightarrow
\text{точка разворота}
\longrightarrow
\text{координата увеличивается}.
\]

На графике \(x(t)\) такой момент часто соответствует \key{излому}.

% ================================================================
\section{3. Скорость по графику \(x(t)\)}

Для прямолинейного равномерного движения скорость определяется по формуле

\[
\boxed{
v=\frac{\Delta x}{\Delta t}
}
\]

где

\[
\Delta x=x_2-x_1,
\qquad
\Delta t=t_2-t_1.
\]

Поэтому

\[
\boxed{
v=\frac{x_2-x_1}{t_2-t_1}
}
\]

Единица измерения скорости:

\[
[v]=\frac{\text{м}}{\text{с}}.
\]

\subsection{Метод 1. «Вертикаль на горизонталь»}

Выберите две удобные точки на одном прямолинейном участке графика.

Затем определите:

\[
\text{изменение координаты по вертикали}
\]

и

\[
\text{изменение времени по горизонтали}.
\]

После этого вычислите

\[
v=
\frac{\text{изменение координаты}}
     {\text{изменение времени}}.
\]

Например:

\[
\Delta x=-6~\text{м},
\qquad
\Delta t=4~\text{с}.
\]

Тогда

\[
v=\frac{-6}{4}=-1{,}5~\frac{\text{м}}{\text{с}}.
\]

Знак «минус» показывает, что координата уменьшается.

\subsection{Метод 2. Формула через две точки}

Пусть

\[
A(t_1;x_1),
\qquad
C(t_2;x_2).
\]

Тогда

\[
v=\frac{x_2-x_1}{t_2-t_1}.
\]

Например,

\[
A(1;4),
\qquad
C(5;-2).
\]

Получаем

\[
v=
\frac{-2-4}{5-1}
=
\frac{-6}{4}
=
-1{,}5~\frac{\text{м}}{\text{с}}.
\]

Оба метода дают один и тот же результат.

% ================================================================
\section{4. Что означает знак скорости}

\begin{center}
\begin{tabularx}{0.92\linewidth}{@{}c X@{}}
\toprule
\textbf{Знак} & \textbf{Физический смысл} \\
\midrule
\(v>0\) &
координата увеличивается; тело движется вдоль положительного направления оси \(Ox\) \\
\(v<0\) &
координата уменьшается; тело движется против положительного направления оси \(Ox\) \\
\(v=0\) &
координата постоянна; тело покоится \\
\bottomrule
\end{tabularx}
\end{center}

\important{Чем круче прямолинейный участок графика}, тем больше модуль скорости.

Если сравниваются два участка, удобно сравнивать величины

\[
\left|\frac{\Delta x}{\Delta t}\right|.
\]

% ================================================================
\section{5. Чек-лист перед ответом}

Перед тем как отвечать на вопрос по графику \(x(t)\), проверьте:

\begin{enumerate}
  \item Что отложено по горизонтальной оси?
  \item Что отложено по вертикальной оси?
  \item Где находятся начальная точка, изломы и пересечения с осью \(t\)?
  \item На каких участках координата увеличивается?
  \item На каких участках координата уменьшается?
  \item Где \(x=0\)?
  \item Где тело меняет направление движения?
  \item На каком участке требуется найти скорость?
  \item Какие две точки этого участка удобно использовать?
  \item Каков знак скорости?
\end{enumerate}

% ================================================================
\section{6. Типичные ошибки}

\begin{itemize}
  \item \key{Путать координату и скорость.}

  Условие \(x=0\) означает нахождение тела в начале координат.

  \[
  x=0 \not\Rightarrow v=0.
  \]

  \item \key{Считать, что наклон графика показывает положение тела.}

  Высота точки графика показывает координату \(x\), а наклон участка связан со скоростью.

  \item \key{Забывать знак скорости.}

  При уменьшении координаты

  \[
  v<0.
  \]

  \item \key{Делить время на координату.}

  Правильная формула:

  \[
  v=\frac{\Delta x}{\Delta t}.
  \]

  \item \key{Брать точки с разных прямолинейных участков.}

  Для скорости на одном участке необходимо выбрать две точки именно этого участка.

  \item \key{Считать точку \(x=0\) остановкой.}

  Пересечение оси \(t\) показывает прохождение через начало координат.
\end{itemize}

% ================================================================
% ТРЕНИРОВОЧНЫЙ БЛОК
% ================================================================
\section{Тренировочный блок}

\textbf{Задание 1.}

Рассмотрите график движения тела.

\begin{center}
\begin{tikzpicture}
\begin{axis}[
  width=11.5cm,
  height=5.7cm,
  axis lines=middle,
  xmin=-0.3,
  xmax=7,
  ymin=-4,
  ymax=5,
  xlabel={$t,\ \text{с}$},
  ylabel={$x,\ \text{м}$},
  xtick={0,1,2,3,4,5,6},
  ytick={-3,-2,-1,0,1,2,3,4},
  grid=both,
  major grid style={draw=gridgray},
  clip=false
]
\addplot[
  accent,
  very thick,
  mark=*,
  mark size=1.8pt
]
coordinates {
  (0,4)
  (4,-2)
  (6,2)
};
\end{axis}
\end{tikzpicture}
\end{center}

Ответьте на вопросы:

\begin{enumerate}[label=\alph*)]
  \item Какова координата тела в момент \(t=0\)?
  \item На каком первом участке координата уменьшается?
  \item В какой момент тело изменяет направление движения?
  \item Какова скорость тела на участке от \(t=4\) с до \(t=6\) с?
\end{enumerate}

\vspace{0.8cm}

\textbf{Задание 2.}

Движение тела задано точками графика

\[
A(1;5),\qquad
B(3;1),\qquad
C(5;-3),\qquad
D(8;3).
\]

Точки \(A,B,C\) лежат на одной прямой, точки \(C,D\) --- на другой.

\begin{enumerate}[label=\alph*)]
  \item Определите скорость на участке \(AC\).
  \item Определите скорость на участке \(CD\).
  \item В какой точке тело изменило направление движения?
  \item Между какими точками тело прошло через начало координат?
\end{enumerate}

\clearpage

\section*{Тренировочный блок --- продолжение}

\textbf{Задание 3.}

График движения состоит из трёх прямолинейных участков:

\[
A(0;-3),\qquad
B(2;1),\qquad
C(5;1),\qquad
D(7;-3).
\]

\begin{enumerate}[label=\alph*)]
  \item Найдите скорость тела на участке \(AB\).
  \item Найдите скорость тела на участке \(BC\).
  \item Найдите скорость тела на участке \(CD\).
  \item На каком участке тело покоилось?
  \item В какие моменты тело проходило через начало координат?
  \item На каком из движущихся участков модуль скорости больше?
\end{enumerate}

\vfill

\begin{center}
\color{darkaccent}
\small
Перед вычислением скорости всегда сначала определяйте,
как меняется координата.
\end{center}

% ================================================================
% ОТВЕТЫ
% ================================================================
\clearpage
\section{Ответы}

\renewcommand{\arraystretch}{1.35}

\begin{center}
\begin{tabularx}{\linewidth}{@{}c X@{}}
\toprule
\textbf{Задание} & \textbf{Ответ} \\
\midrule

1 &
а) \(4\) м;
б) от \(t=0\) до \(t=4\) с;
в) \(t=4\) с;
г)
\(
v=\dfrac{2-(-2)}{6-4}=2~\text{м/с}.
\)
\\

\addlinespace

2 &
а)
\(
v_{AC}=\dfrac{-3-5}{5-1}=-2~\text{м/с};
\)
\newline
б)
\(
v_{CD}=\dfrac{3-(-3)}{8-5}=2~\text{м/с};
\)
\newline
в) в точке \(C\);
\newline
г) на участках \(BC\) и \(CD\).
\\

\addlinespace

3 &
а)
\(
v_{AB}=\dfrac{1-(-3)}{2-0}=2~\text{м/с};
\)
\newline
б)
\(
v_{BC}=0;
\)
\newline
в)
\(
v_{CD}=\dfrac{-3-1}{7-5}=-2~\text{м/с};
\)
\newline
г) \(BC\);
\newline
д)
\(
t=1{,}5~\text{с}
\)
и
\(
t=5{,}5~\text{с};
\)
\newline
е) модули скоростей на участках \(AB\) и \(CD\) одинаковы:
\(
|v|=2~\text{м/с}.
\)
\\

\bottomrule
\end{tabularx}
\end{center}

\vfill

\begin{center}
{\large\bfseries\color{darkaccent}
Главная идея}
\end{center}

\[
\boxed{
\text{сначала анализируем изменение координаты, затем вычисляем скорость}
}
\]

\[
\boxed{
v=\frac{\Delta x}{\Delta t}
}
\]

\begin{center}
\small
Если координата растёт, скорость положительна.\\
Если координата уменьшается, скорость отрицательна.\\
Если координата постоянна, скорость равна нулю.
\end{center}

\end{document}